% GENERATED FILE. Edit canonical lessons, then run npm run generate:books.
% canonical-source-hash: fnv1a64:f34c02b1c58ced01
% canonical-lessons: KA-C386-samajayisi, KA-C386-sabubu, KA-C386-akasmika, KA-C386-bimbisu, KA-C386-hone

\chapter{Excuses and Explanations}
\label{ch:ka-excuses-and-explanations}
\hlchaptermodality{386}

\begin{chapteropening}
\textbf{By the end of this chapter:} \emph{I can say an explanation, a lame excuse, accidental, to portray and blame.}

Complete the last lesson of chapter 386: I can say an explanation, a lame excuse, accidental, to portray and blame.
\end{chapteropening}

\section[\texorpdfstring{\kn{ಸಮಜಾಯಿಷಿ} (samajāyiṣi)}{ಸಮಜಾಯಿಷಿ (samajāyiṣi)}]{\kn{ಸಮಜಾಯಿಷಿ} (samajāyiṣi) — an explanation}

\label{lesson:KA-C386-samajayisi}



\begin{quote}
Before the new one: say the Kannada for a tendency, then the Kannada for a source.
\end{quote}

\subsection*{You'll want to know: \kn{ಸಮಜಾಯಿಷಿ}}
\textbf{\kn{ಸಮಜಾಯಿಷಿ}} — \emph{samajāyiṣi} — \textquotedblleft{}an explanation\textquotedblright{}.

\begin{quote}
\kn{ತಡವಾಗಿ} \kn{ಬಂದಿದ್ದಕ್ಕೆ} \kn{ಅವನು} \kn{ಸಮಜಾಯಿಷಿ} \kn{ಕೊಟ್ಟನು}.
\end{quote}

\begin{quote}
\emph{taḍavāgi bandiddakke avanu samajāyiṣi koṭṭanu.}
\end{quote}

\begin{quote}
He gave an explanation for coming late.
\end{quote}

\kn{ಸಮಜಾಯಿಷಿ} \kn{ಕೊಡು}, to explain oneself.

\begin{cousinweb}
\textbf{\kn{ಸಮಜಾಯಿಷಿ}} (\emph{samajāyiṣi}) is usually traced to Hindustani \emph{samjhāish}, \textquotedblleft{}an explanation\textquotedblright{}, from \emph{samajhnā}, \textquotedblleft{}to understand\textquotedblright{}; DEDR and CDIAL do not list the Kannada form.
\end{cousinweb}

\begin{grammarlens}[title={What you've built}]
One more word for giving reasons and opinions.
\end{grammarlens}

\subsection*{Your turn}
\begin{itemize}
\raggedright
  \item \emph{Say it:} \emph{samajāyiṣi}
  \item \emph{Say it:} \emph{samajāyiṣi}, once more
  \item \emph{Say it:} \emph{taḍavāgi bandiddakke avanu samajāyiṣi koṭṭanu.} — He gave an explanation for coming late.
  \item \emph{Recall:} say the Kannada for a tendency, then the Kannada for a source, then say \emph{samajāyiṣi} again
\end{itemize}

\begin{tcolorbox}[breakable,colback=teal!4,colframe=teal!35!black,title={Before you move on}]
What does \kn{ಸಮಜಾಯಿಷಿ} mean? (\textquotedblleft{}An explanation\textquotedblright{}.) Say it once more.
\end{tcolorbox}
\section[\texorpdfstring{\kn{ಸಬೂಬು} (sabūbu)}{ಸಬೂಬು (sabūbu)}]{\kn{ಸಬೂಬು} (sabūbu) — a lame excuse}

\label{lesson:KA-C386-sabubu}



\begin{quote}
Before the new one: say the Kannada for a source, then the Kannada for an explanation.
\end{quote}

\subsection*{You'll want to know: \kn{ಸಬೂಬು}}
\textbf{\kn{ಸಬೂಬು}} — \emph{sabūbu} — \textquotedblleft{}a lame excuse\textquotedblright{}.

\begin{quote}
\kn{ಕೆಲಸ} \kn{ಮಾಡದಿರಲು} \kn{ಅವನು} \kn{ಸಬೂಬು} \kn{ಹೇಳಿದನು}.
\end{quote}

\begin{quote}
\emph{kelasa māḍadiralu avanu sabūbu hēḷidanu.}
\end{quote}

\begin{quote}
He made an excuse for not working.
\end{quote}

\kn{ಸಬೂಬು} \kn{ಹೇಳು}, to make excuses.

\begin{cousinweb}
\textbf{\kn{ಸಬೂಬು}} (\emph{sabūbu}) is usually traced to Arabic \emph{sabab}, \textquotedblleft{}a cause, a reason\textquotedblright{}, which reached Kannada through Hindustani.
\end{cousinweb}

\begin{grammarlens}[title={What you've built}]
One more word for giving reasons and opinions.
\end{grammarlens}

\subsection*{Your turn}
\begin{itemize}
\raggedright
  \item \emph{Say it:} \emph{sabūbu}
  \item \emph{Say it:} \emph{sabūbu}, once more
  \item \emph{Say it:} \emph{kelasa māḍadiralu avanu sabūbu hēḷidanu.} — He made an excuse for not working.
  \item \emph{Recall:} say the Kannada for a source, then the Kannada for an explanation, then say \emph{sabūbu} again
\end{itemize}

\begin{tcolorbox}[breakable,colback=teal!4,colframe=teal!35!black,title={Before you move on}]
What does \kn{ಸಬೂಬು} mean? (\textquotedblleft{}A lame excuse\textquotedblright{}.) Say it once more.
\end{tcolorbox}
\section[\texorpdfstring{\kn{ಆಕಸ್ಮಿಕ} (ākasmika)}{ಆಕಸ್ಮಿಕ (ākasmika)}]{\kn{ಆಕಸ್ಮಿಕ} (ākasmika) — accidental; sudden}

\label{lesson:KA-C386-akasmika}



\begin{quote}
Before the new one: say the Kannada for an explanation, then the Kannada for a lame excuse.
\end{quote}

\subsection*{You'll want to know: \kn{ಆಕಸ್ಮಿಕ}}
\textbf{\kn{ಆಕಸ್ಮಿಕ}} — \emph{ākasmika} — \textquotedblleft{}accidental; sudden\textquotedblright{}.

\begin{quote}
\kn{ನಮ್ಮ} \kn{ಭೇಟಿ} \kn{ಆಕಸ್ಮಿಕವಾಗಿತ್ತು}.
\end{quote}

\begin{quote}
\emph{namma bhēṭi ākasmikavāgittu.}
\end{quote}

\begin{quote}
Our meeting was accidental.
\end{quote}

\kn{ಆಕಸ್ಮಿಕವಾಗಿ}, by chance.

\begin{cousinweb}
\textbf{\kn{ಆಕಸ್ಮಿಕ}} (\emph{ākasmika}) is borrowed from Sanskrit \emph{ākasmika}, \textquotedblleft{}causeless, unforeseen\textquotedblright{} (Monier-Williams), from \emph{akasmāt}, \textquotedblleft{}without a why\textquotedblright{}.
\end{cousinweb}

\begin{grammarlens}[title={What you've built}]
One more word for giving reasons and opinions.
\end{grammarlens}

\subsection*{Your turn}
\begin{itemize}
\raggedright
  \item \emph{Say it:} \emph{ākasmika}
  \item \emph{Say it:} \emph{ākasmika}, once more
  \item \emph{Say it:} \emph{namma bhēṭi ākasmikavāgittu.} — Our meeting was accidental.
  \item \emph{Recall:} say the Kannada for an explanation, then the Kannada for a lame excuse, then say \emph{ākasmika} again
\end{itemize}

\begin{tcolorbox}[breakable,colback=teal!4,colframe=teal!35!black,title={Before you move on}]
What does \kn{ಆಕಸ್ಮಿಕ} mean? (\textquotedblleft{}Accidental; sudden\textquotedblright{}.) Say it once more.
\end{tcolorbox}
\section[\texorpdfstring{\kn{ಬಿಂಬಿಸು} (bimbisu)}{ಬಿಂಬಿಸು (bimbisu)}]{\kn{ಬಿಂಬಿಸು} (bimbisu) — to portray; to present as}

\label{lesson:KA-C386-bimbisu}



\begin{quote}
Before the new one: say the Kannada for a lame excuse, then the Kannada for accidental.
\end{quote}

\subsection*{You'll want to know: \kn{ಬಿಂಬಿಸು}}
\textbf{\kn{ಬಿಂಬಿಸು}} — \emph{bimbisu} — \textquotedblleft{}to portray; to present as\textquotedblright{}.

\begin{quote}
\kn{ಅವನು} \kn{ತನ್ನನ್ನು} \kn{ಒಳ್ಳೆಯವನಂತೆ} \kn{ಬಿಂಬಿಸಿದನು}.
\end{quote}

\begin{quote}
\emph{avanu tannannu oḷḷeyavanante bimbisidanu.}
\end{quote}

\begin{quote}
He made himself out to be a good man.
\end{quote}

Past (I …): \kn{ಬಿಂಬಿಸಿದೆ}, \emph{bimbiside}.

As what: -\kn{ಅಂತೆ}, like. \kn{ಬಿಂಬ}, an image.

\begin{cousinweb}
\textbf{\kn{ಬಿಂಬಿಸು}} (\emph{bimbisu}) is built on Sanskrit \emph{bimba}, \textquotedblleft{}a round form; a form and its reflection, an image\textquotedblright{} (CDIAL 9242), with the Kannada verb ending -\emph{isu}.
\end{cousinweb}

\begin{grammarlens}[title={What you've built}]
One more word for giving reasons and opinions.
\end{grammarlens}

\subsection*{Your turn}
\begin{itemize}
\raggedright
  \item \emph{Say it:} \emph{bimbisu}
  \item \emph{Say it:} \emph{bimbisu}, once more
  \item \emph{Say it:} \emph{avanu tannannu oḷḷeyavanante bimbisidanu.} — He made himself out to be a good man.
  \item \emph{Recall:} say the Kannada for a lame excuse, then the Kannada for accidental, then say \emph{bimbisu} again
\end{itemize}

\begin{tcolorbox}[breakable,colback=teal!4,colframe=teal!35!black,title={Before you move on}]
What does \kn{ಬಿಂಬಿಸು} mean? (\textquotedblleft{}To portray; to present as\textquotedblright{}.) Say it once more.
\end{tcolorbox}
\section[\texorpdfstring{\kn{ಹೊಣೆ} (hoṇe)}{ಹೊಣೆ (hoṇe)}]{\kn{ಹೊಣೆ} (hoṇe) — blame; responsibility}

\label{lesson:KA-C386-hone}



\begin{quote}
Before the new one: say the Kannada for accidental, then the Kannada for to portray.
\end{quote}

\subsection*{You'll want to know: \kn{ಹೊಣೆ}}
\textbf{\kn{ಹೊಣೆ}} — \emph{hoṇe} — \textquotedblleft{}blame; responsibility\textquotedblright{}.

\begin{quote}
\kn{ಈ} \kn{ನಷ್ಟಕ್ಕೆ} \kn{ಯಾರು} \kn{ಹೊಣೆ}?
\end{quote}

\begin{quote}
\emph{ī naṣṭakke yāru hoṇe?}
\end{quote}

\begin{quote}
Who is to blame for this loss?
\end{quote}

\kn{ಹೊಣೆ} \kn{ಹೊರು}, to take the blame.

\begin{cousinweb}
\textbf{\kn{ಹೊಣೆ}} (\emph{hoṇe}) is an everyday Kannada word whose origin is not settled: neither DEDR nor CDIAL lists it.
\end{cousinweb}

\begin{grammarlens}[title={What you've built}]
One more word for giving reasons and opinions.
\end{grammarlens}

\subsection*{Your turn}
\begin{itemize}
\raggedright
  \item \emph{Say it:} \emph{hoṇe}
  \item \emph{Say it:} \emph{hoṇe}, once more
  \item \emph{Say it:} \emph{ī naṣṭakke yāru hoṇe?} — Who is to blame for this loss?
  \item \emph{Recall:} say the Kannada for accidental, then the Kannada for to portray, then say \emph{hoṇe} again
\end{itemize}

\begin{tcolorbox}[breakable,colback=teal!4,colframe=teal!35!black,title={Before you move on}]
What does \kn{ಹೊಣೆ} mean? (\textquotedblleft{}Blame; responsibility\textquotedblright{}.) Say it once more.
\end{tcolorbox}
